\section*{Chapitre \ref{ch:b2:magma-emplacement} --- Magmas et mise en place}

\begin{solution}{exo:b2:magma-emplacement:1}
$T = \qty{1073}{K}$. Pour $w = 1$, $\ln w = 0$~: $\log_{10}\eta = -3.545 +
9601/(1073 - 195.7) = -3.545 + 10.94 = 7.40$. Pour $w = 4$, $\ln 4 = 1.386$~:
$A = -3.545 + 1.155 = -2.390$, $B = 9601 - 3283 = 6318$, $C = 195.7 + 44.7 =
240.4$, et $\log_{10}\eta = -2.390 + 6318/832.6 = -2.390 + 7.588 = 5.20$. La
viscosité diminue d'un facteur $10^{7.40 - 5.20} = 10^{2.20} \approx 160$.
\end{solution}

\begin{solution}{exo:b2:magma-emplacement:2}
Rhyolite~: $p_s = (6/0.44)^2 = \qty{186}{MPa}$, à une profondeur de $\num{186e6}/(2600\times9.81)
= \qty{7.3}{km}$. Basalte~: $p_s = (0.5/0.34)^2 = \qty{2.16}{MPa}$, à
$\num{2.16e6}/(2600\times9.81) = \qty{85}{m}$. La rhyolite commence à former des bulles
en profondeur dans la croûte, le basalte seulement près de la surface (s'il ne contient pas de CO$_2$).
\end{solution}

\begin{solution}{exo:b2:magma-emplacement:3}
$(1150 + 50)/2 = \qty{600}{\celsius}$ et $(1150 + 400)/2 = \qty{775}{\celsius}$. Le
\omterm{def:b2:magma-emplacement:bodies}{dyke} dans la roche froide est refroidi plus fortement à son contact, et a la \omterm{def:b2:magma-emplacement:aureole}{bordure figée}
la plus large.
\end{solution}

\begin{solution}{exo:b2:magma-emplacement:4}
(a) Un \omterm{def:b2:magma-emplacement:bodies}{sill}, parallèle à la stratification. (b) Un \omterm{def:b2:magma-emplacement:bodies}{dyke}, recoupant la stratification.
(c) Un \omterm{def:b2:magma-emplacement:bodies}{batholite}, de plus de \qty{100}{km^2}, assemblé à partir de nombreuses intrusions.
(d) Une \omterm{def:b2:magma-emplacement:bodies}{laccolite}, qui a soulevé son toit en dôme. (e) Un \omterm{def:b2:magma-emplacement:bodies}{dyke}~: le granite n'a pas de stratification,
de sorte que toute nappe qui le traverse est un \omterm{def:b2:magma-emplacement:bodies}{dyke}, quelle que soit son orientation.
\end{solution}

\begin{solution}{exo:b2:magma-emplacement:5}
$(1 - 1.35\phi)^{-2.5}$~: pour $\phi = 0.05$, $0.9325^{-2.5} = 1.19$, soit
\qty{119}{Pa.s}, contre $1 + 2.5\times0.05 = 1.125$, \qty{113}{Pa.s}, avec la formule d'Einstein
(proche, comme il se doit à la limite diluée). Pour $0.20$~: $0.73^{-2.5} =
2.20$, \qty{220}{Pa.s} ; pour $0.40$~: $0.46^{-2.5} = 6.97$, \qty{697}{Pa.s} ; pour $0.60$~:
$0.19^{-2.5} = 63.6$, soit environ \qty{6360}{Pa.s}. Blocage à $\phi = 1/1.35 = 0.74$.
\end{solution}

\begin{solution}{exo:b2:magma-emplacement:6}
$G = 300\times9.81 = \qty{2943}{Pa/m}$. (a) $\bar u = 2943\times1^2/(12\times100) =
\qty{2.45}{m/s}$ ; $q = \qty{2.45}{m^2/s}$ ; $Re = 2700\times2.45\times1/100 = 66$, laminaire ;
$t = 10000/2.45 = \qty{4080}{s}$, soit environ \qty{1.1}{h}. (b) $\bar u =
2943/(12\times\num{e6}) = \qty{2.45e-4}{m/s}$, $t = \qty{4.08e7}{s} = 472$ jours ; avec
$w = \qty{10}{m}$, $\bar u$ est cent fois plus grande, \qty{0.0245}{m/s}, et $t =
\qty{4.08e5}{s} = 4.7$ jours. (c) Le \omterm{def:b2:magma-emplacement:bodies}{dyke} de \qty{1}{m} se refroidit en environ
$0.5^2/\num{e-6} = \qty{2.5e5}{s}$, soit 2.9 jours ; le \omterm{def:b2:magma-emplacement:bodies}{dyke} de \qty{10}{m} en $5^2/\num{e-6} =
\qty{2.5e7}{s}$, soit 289 jours. Le basalte (\qty{1.1}{h}) et la rhyolite dans le \omterm{def:b2:magma-emplacement:bodies}{dyke} large
(4.7 jours) arrivent avant que leur \omterm{def:b2:magma-emplacement:bodies}{dyke} ne gèle ; la rhyolite dans le
\omterm{def:b2:magma-emplacement:bodies}{dyke} de \qty{1}{m} (472 jours) ne le peut pas.
\end{solution}

\begin{solution}{exo:b2:magma-emplacement:7}
Par kilogramme de \omterm{def:b2:magma-emplacement:magma}{magma}, $n = 10/18.015 = \qty{0.555}{mol}$ d'eau et $nRT =
0.555\times8.314\times1173 = \qty{5414}{J}$ ; le liquide restant occupe $0.99/2400 =
\qty{4.13e-4}{m^3}$. À \qty{100}{MPa} : $V = \qty{5.41e-5}{m^3}$, fraction
$5.41/(5.41 + 41.3) = 0.12$. À \qty{10}{MPa} : $V = \qty{5.41e-4}{m^3}$, fraction
$0.57$. À \qty{0.1}{MPa} : $V = \qty{0.0541}{m^3}$, fraction $0.992$. La valeur à
\qty{100}{MPa} est la moins fiable : la vapeur d'eau à une telle pression est dense et
loin d'un gaz parfait (et une grande partie de l'eau resterait en fait dissoute).
\end{solution}

\begin{solution}{exo:b2:magma-emplacement:8}
(a) $2300 + 60z = 2650$ donne $z = \qty{5.8}{km}$. (b) Le liquide a
\qty{2580}{kg/m^3} : $z = 280/60 = \qty{4.7}{km}$, soit environ \qty{1.2}{km} plus haut. (c)
À $z = \qty{7.8}{km}$, la roche a $2300 + 60\times7.8 = \qty{2770}{kg/m^3}$ ; la
force vaut $(2770 - 2650)\times9.81 = \qty{1.2}{kN}$ par mètre cube, dirigée vers le haut.
\end{solution}

\begin{solution}{exo:b2:magma-emplacement:9}
$v = 2\times300\times9.81\times(1000)^2/(9\times\num{e19}) = \qty{6.5e-11}{m/s}$, soit
$\num{6.5e-11}\times\num{3.156e7} = \qty{2.1e-3}{m}$ par an, \qty{2.1}{mm/an}. La remontée de
\qty{10}{km} prend $\num{e4}/\num{2.1e-3} = \qty{4.8}{Ma}$. L'échelle de temps du refroidissement
est $(1000)^2/\num{e-6} = \qty{e12}{s} = \qty{32}{ka}$ : le corps se solidifie après
avoir monté d'environ $32000\times0.0021 \approx \qty{70}{m}$. Un \omterm{def:b2:magma-emplacement:ascent}{diapir} magmatique ne peut pas traverser
une telle croûte ; il lui faut un environnement plus chaud et plus mou.
\end{solution}

\begin{solution}{exo:b2:magma-emplacement:10}
(a) $x_{\mathrm{OH}} = \sqrt{K_2x_m} = \sqrt{K_1K_2f}$, donc $x = K_1f + \frac12\sqrt{K_1K_2f}$.
Le rapport du premier terme au second vaut $2\sqrt{K_1f/K_2}$ : il tend vers zéro
lorsque $f \to 0$ (la racine carrée domine) et croît sans borne avec $f$ (le
terme linéaire domine). (b) $K_1 f = \frac12\sqrt{K_1K_2 f}$ donne $K_1^2f^2 =
K_1K_2f/4$, donc $f = K_2/(4K_1)$. (c) $f = 0.2/(\num{4e-4}) = \qty{500}{MPa}$. À
\qty{100}{MPa} : $x_m = 0.0100$, $\frac12\sqrt{\num{e-4}\times0.2\times100} =
\frac12\times0.0447 = 0.0224$, donc $x = 0.0324$ et la part hydroxyle vaut
$0.0224/0.0324 = \qty{69}{\percent}$.
\end{solution}

\begin{solution}{exo:b2:magma-emplacement:11}
(a) $1000/\num{4e6} = \qty{2.5e-4}{km^3}$ par an, soit $\qty{2.5e5}{m^3}$ par an,
ou $\num{2.5e5}/\num{3.156e7} = \qty{0.0079}{m^3/s}$, soit environ huit litres par seconde.
(b) Le \omterm{def:b2:magma-emplacement:bodies}{dyke} transporte $1000\times2\times0.8 = \qty{1600}{m^3/s}$ et apporte
$\qty{e12}{m^3}$ en \qty{6.25e8}{s}, soit environ 20 ans : $\num{5e-6}$ de la
\qty{4}{Ma}. (c) Chaque nappe contient $20\times\num{e4}\times\num{e4} = \qty{2e9}{m^3} =
\qty{2}{km^3}$, soit 500 nappes, une toutes les $\num{4e6}/500 = 8000$ ans en moyenne.
(d) $t = (10)^2/(4\times0.65^2\times\num{0.7e-6}) = \qty{8.5e7}{s}$, soit 2.7 ans. Chaque
nappe se solidifie des milliers d'années avant l'arrivée de la suivante : sauf si l'apport
arrive par à-coups rapides, le \omterm{def:b2:magma-emplacement:bodies}{pluton} n'est jamais un seul grand corps de \omterm{def:b2:magma-emplacement:magma}{magma}, comme le
suggère la datation des grands \omterm{def:b2:magma-emplacement:bodies}{plutons}.
\end{solution}

\begin{solution}{exo:b2:magma-emplacement:12}
(a) À $x$ fixé, $\ln(T - T_0) = \text{cste} - \frac12\ln t - x^2/(4\kappa t)$, dont
dérivée $-1/(2t) + x^2/(4\kappa t^2)$ s'annule en $t^* = x^2/(2\kappa)$. Alors
$e^{-x^2/(4\kappa t^*)} = e^{-1/2}$ et $\sqrt{4\pi\kappa t^*} = x\sqrt{2\pi}$, donc $\Delta
T_{\max} = (T_m - T_0)\,b\,e^{-1/2}/(x\sqrt{2\pi}) = (T_m - T_0)\,b/(x\sqrt{2\pi e})$. (b)
$\sqrt{2\pi e} = 4.13$ : $x = 1100\times2/(100\times4.13) = \qty{5.3}{m}$, atteint après
$t^* = 5.3^2/(2\times\num{0.7e-6}) = \qty{2.0e7}{s}$, soit environ 230 jours (à seulement 2.7
largeurs de \omterm{def:b2:magma-emplacement:bodies}{dyke}, l'approximation de source plane est grossière). (c) Le réchauffement décroît
comme $b/x$ : un \omterm{def:b2:magma-emplacement:bodies}{dyke} mince transporte peu de chaleur, et la roche ne se réchauffe sensiblement
que dans un rayon de quelques largeurs de \omterm{def:b2:magma-emplacement:bodies}{dyke}.
\end{solution}

\begin{solution}{pb:b2:magma-emplacement:1}
\textbf{1.} $\log_{10}\eta = -4.55 + 5510/(1423 - 611) = -4.55 + 6.786 = 2.24$~;
$\eta = \qty{172}{Pa.s}$.

\textbf{2.} $-4.55 + 5510/762 = -4.55 + 7.231 = 2.68$, $\eta = \qty{480}{Pa.s}$~: un
facteur de $480/172 = 2.8$ pour \qty{50}{K}.

\textbf{3.} $172\times(1 - 1.35\times0.20)^{-2.5} = 172\times0.73^{-2.5} = 172\times2.20 =
\qty{378}{Pa.s}$.

\textbf{4.} $\phi = 1/1.35 = 0.74$. Les observations de terrain placent la fin de l'écoulement d'ensemble
à un taux de fusion de \qtyrange{30}{50}{\percent}, soit
\qtyrange{50}{70}{\percent} de cristaux : plus tôt, parce que les cristaux réels ne sont pas
des sphères d'une seule taille.

\textbf{5.} Socle~: $2800 - 2710 = \qty{90}{kg/m^3}$, le \omterm{def:b2:magma-emplacement:magma}{magma} est plus léger.
Sédiments~: $2450 - 2710 = \qty{-260}{kg/m^3}$, le \omterm{def:b2:magma-emplacement:magma}{magma} est plus dense.

\textbf{6.} $p = 2450\times9.81\times3000 = \qty{72.1}{MPa}$.

\textbf{7.} $90\times9.81 = \qty{883}{Pa/m}$.

\textbf{8.} $\bar u = 883\times2.0^2/(12\times378) = 3532/4536 = \qty{0.78}{m/s}$.

\textbf{9.} $q = 0.78\times2.0 = \qty{1.56}{m^2/s}$ ; sur \qty{1.0}{km}, soit environ
\qty{1560}{m^3/s}.

\textbf{10.} $Re = 2710\times0.78\times2.0/378 = 11$~: laminaire.

\textbf{11.} $12000/0.78 = \qty{15400}{s}$, soit environ \qty{4.3}{h}.

\textbf{12.} $(1.0)^2/\num{0.7e-6} = \qty{1.43e6}{s}$, soit 16.5 jours, bien plus long que
\qty{4.3}{h} : le \omterm{def:b2:magma-emplacement:bodies}{dyke} reste ouvert pendant que le \omterm{def:b2:magma-emplacement:magma}{magma} s'écoule.

\textbf{13.} À \qty{3.0}{km}, le \omterm{def:b2:magma-emplacement:magma}{magma} rencontre des sédiments plus légers que lui : la
base des sédiments est son \omterm{def:b2:magma-emplacement:ascent}{niveau de flottabilité neutre}. Au-dessus, le \omterm{def:b2:magma-emplacement:magma}{magma} serait
plus dense que son environnement, il s'étale donc latéralement le long du contact, en
\omterm{def:b2:magma-emplacement:bodies}{sill}.

\textbf{14.} $p_s = (0.40/0.34)^2 = \qty{1.38}{MPa}$, atteinte à
$\num{1.38e6}/(2450\times9.81) = \qty{58}{m}$ sous la surface.

\textbf{15.} $0.34\times\sqrt{72.1} = 0.34\times8.49 = \qty{2.9}{\percent}$, bien plus que
les \qty{0.40}{\percent} présents : le \omterm{def:b2:magma-emplacement:magma}{magma} du \omterm{def:b2:magma-emplacement:bodies}{sill} n'exsolve pas d'eau.

\textbf{16.} Le CO$_2$ est beaucoup moins soluble que l'eau : il sature et forme des bulles
en profondeur pendant la remontée, si bien que le premier gaz d'un basalte est surtout du CO$_2$.

\textbf{17.} Par kilogramme, $n = 4.0/18.015 = \qty{0.222}{mol}$, $V = 0.222\times8.314\times1423/
\num{e5} = \qty{0.0263}{m^3}$ ; le liquide occupe $0.996/2710 = \qty{3.68e-4}{m^3}$ ; la
fraction vaut $0.0263/(0.0263 + 0.000368) = 0.986$, soit \qty{98.6}{\percent}.

\textbf{18.} $(1150 + 100)/2 = \qty{625}{\celsius}$.

\textbf{19.} Après 10 ans, $2\sqrt{\kappa t} = 2\sqrt{\num{0.7e-6}\times\num{3.156e8}} =
\qty{29.7}{m}$. Le point est à $x = \qty{60}{m}$ du milieu : $T = 100 +
525\,[\operatorname{erf}(-10/29.7) + \operatorname{erf}(110/29.7)] = 100 + 525\,(-0.365 +
1.000) = \qty{433}{\celsius}$.

\textbf{20.} Au milieu, $T = 100 + 1050\operatorname{erf}(50/2\sqrt{\kappa t})$ : après
10 ans $1050\operatorname{erf}(1.68) = 1050\times0.983$, $T = \qty{1132}{\celsius}$ ; après
50 ans, $2\sqrt{\kappa t} = \qty{66.5}{m}$, $1050\operatorname{erf}(0.752) = 1050\times0.713$,
$T = \qty{848}{\celsius}$.

\textbf{21.} À $x = \qty{130}{m}$ : \qty{294}{\celsius} après 300 ans,
\qty{295}{\celsius} après 363 ans (le maximum), \qty{295}{\celsius} après 420
ans. Sans chaleur latente, la roche située à \qty{80}{m} au-dessus du \omterm{def:b2:magma-emplacement:bodies}{sill} n'atteint jamais
tout à fait \qty{300}{\celsius}.

\textbf{22.} $\num{9e8}\times1.772/(\num{2.5e6}\times1050) = 0.608$.

\textbf{23.} $\lambda = 0.64$~: $e^{-0.4096}/(0.64\times1.6346) = 0.6639/1.0461 = 0.635$ ;
$0.65$~: $0.6554/1.0673 = 0.614$ ; $0.66$~: $0.6469/1.0886 = 0.594$. La valeur $0.608$
se situe entre $0.65$ et $0.66$, proche de $0.65$~: $\lambda = 0.65$.

\textbf{24.} $100 + 1050/(1 + 0.642) = 100 + 640 = \qty{740}{\celsius}$. La chaleur latente
libérée au front de solidification maintient le côté \omterm{def:b2:magma-emplacement:magma}{magma} à
\qty{1150}{\celsius} plus longtemps, si bien que le contact reçoit davantage de chaleur.

\textbf{25.} $(50)^2/\num{0.7e-6} = \qty{3.57e9}{s}$, soit 113 ans.

\textbf{26.} $t = (100)^2/(16\times0.65^2\times\num{0.7e-6}) = \num{e4}/\num{4.73e-6} =
\qty{2.11e9}{s}$, soit $\num{2.11e9}/\num{3.156e7}$ : le \omterm{def:b2:magma-emplacement:bodies}{sill} est solide après
\textbf{\qty[text-series-to-math=true]{67}{} ans}. Il est plus court que
$(b/2)^2/\kappa$, d'un facteur $1/(4\lambda^2) = 0.59$ : cette échelle de temps ne fixe que
l'ordre de grandeur, et les deux parois du \omterm{def:b2:magma-emplacement:bodies}{sill} extraient la chaleur d'un encaissant froid à
la fois, assez vite pour évacuer aussi la chaleur latente.
\end{solution}
